% ============================================================
% Chapter 4: Proposed Method: Specification and Design
%   Section: Requirements
%     Subsection: Non-Functional Requirements
% Included from chapters/04-proposed-method/01-requirements/requirements.tex
% ============================================================

\subsection{Non-Functional Requirements}

\begin{table}[ht]
\centering
\small
\caption{Non-functional requirements and how each is met.}
\label{tab:nonfunctional-requirements}
\begin{tabularx}{\linewidth}{@{}lXX@{}}
\toprule
\textbf{Id} & \textbf{Requirement} & \textbf{How it is met} \\
\midrule
NFR1 & \textbf{Low cost}: no paid API; runs on a laptop or one VM &
Local open-weight models; self-hosted metasearch; one compose stack \\
NFR2 & \textbf{Auditability}: every verdict traceable to its evidence &
Cited sources with verbatim quotes; a per-run journal of every event;
lineage on every stored record \\
NFR3 & \textbf{Fail towards caution} & \texttt{UNVERIFIED} when the
evidence does not decide; ungrounded verdicts downgraded; outages
reported apart \\
NFR4 & \textbf{Two languages} & English and Spanish lexicons,
language detection, per-language queries; a third language fails the
test suite \\
NFR5 & \textbf{Reproducible evaluation} & Resumable runs keyed by
dataset, model, thresholds and commit; seeded bootstrap; committed
reports \\
NFR6 & \textbf{Per-run tunability} & Every threshold can be overridden
for one run without affecting concurrent runs \\
NFR7 & \textbf{Bounded load on shared services} & One concurrency
ceiling per external service, not per request \\
NFR8 & \textbf{Security} & An address-based guard on every server-side
fetch; an API key and a site password in production; read-only graph
console \\
\bottomrule
\end{tabularx}
\end{table}
